\section{三个 LLM moments：目标函数、规模与对齐}

讲者用 BERT、GPT-3、ChatGPT 三个时刻压缩五年语言模型史。它们分别改变了研究者对 self-supervision、scale 和 task adaptation 的理解。关键不是记住年份，而是看研究瓶颈如何从“设计更好的任务”迁移到“分配计算”，再迁移到“用少量高质量数据改变行为”。

\subsection{BERT moment：理解与生成曾被认为需要不同目标}

2018--2021 年常见分工是：masked language modeling（MLM）负责理解，autoregressive modeling 负责生成，T5 用 encoder--decoder 兼顾两者。GLM 的 blank infilling 则把被遮盖 span 内部改成自回归生成，试图在一个 decoder-style objective 中覆盖不同任务。

\lecturefigure{slide-03-language-modeling.jpg}{BERT 时刻：MLM、autoregression、T5 与 GLM 的任务视角}{Stanford Online 官方录像 00:07:00}

设原序列为 $x=(x_1,\ldots,x_n)$，遮盖 span 集合为 $M$。GLM 的核心可写成对被遮盖 token 的条件自回归：
\[
\mathcal{L}_{\text{GLM}}
=-\sum_{m\in M}\sum_{j=1}^{|m|}
\log p_\theta\!\left(x_{m,j}\mid x_{\setminus M},x_{m,<j}\right).
\]
这里 $x_{\setminus M}$ 是可见上下文，$x_{m,<j}$ 是当前 span 内已生成 token。若遮盖几乎覆盖全序列，目标接近 GPT；若遮盖较小 span，则保留双向可见上下文。

\begin{knowledgebox}{Objective 的真正作用}
预训练 objective 决定模型在训练时能看到什么条件、预测什么变量以及误差怎样传播。它不直接等于“理解”或“生成”能力；下游 fine-tuning、prompt format、data mixture 和 inference procedure 都会改变最终行为。
\end{knowledgebox}

\teachervoice{00:05:00--00:08:55，讲者把 GLM 放在“社区信念会快速改变”的历史里：当年大家努力发明 self-supervised task，后来发现 task adaptation 可以很便宜，研究价值的判断随规模和数据条件变化。}

\subsection{GPT-3 moment：scaling law 变成预算分配器}

第二个时刻不是“参数越多越好”的口号，而是用可拟合曲线把 loss 与参数、数据、计算联系起来。常见经验形式为
\[
L(N,D)\approx L_\infty+aN^{-\alpha}+bD^{-\beta},
\]
其中 $N$ 是参数量，$D$ 是训练 token 数，$L_\infty$ 是不可约 loss，$a,b,\alpha,\beta$ 由实验拟合。它允许团队在固定 compute $C$ 下比较“加参数”与“加 token”的边际收益。

\lecturefigure{slide-04-scaling-law.jpg}{GPT-3 时刻：用 scaling curves 分配参数、数据与计算}{Stanford Online 官方录像 00:09:19}

若一次 dense Transformer 训练的主要 FLOPs 近似为
\[
C\approx 6ND,
\]
那么给定预算时，$N$ 与 $D$ 不能同时无限增长。这里的 $6$ 是常用粗略常数，真实成本还受 sequence length、attention、activation recomputation、parallel efficiency 和 kernels 影响。

\begin{warningbox}{Scaling law 不是性能保证书}
拟合曲线只在观测到的模型族、数据分布、tokenizer、optimizer 和训练区间内成立。数据质量下降、重复率上升、分布外任务或系统效率恶化时，增加名义 compute 不保证按同一指数提升。
\end{warningbox}

\teachervoice{00:09:19--00:10:00，讲者把 scaling law 形容成工程预算工具：老板给四倍资源时，需要决定资源落在参数还是 token，而不是只宣布“模型更大”。}

\subsection{ChatGPT moment：task adaptation 便宜，但知识从哪里来？}

第三张 slide 把 InstructGPT 的人类偏好提升与“下游性能由 pretraining loss 决定”的经验图放在一起。前者说明 alignment 能显著改变人类评价；后者试图说明同等预训练 loss 下，较小但训练更充分的模型可能追上较大但 under-trained 的模型。

\lecturefigure{slide-05-alignment-pretraining-loss.jpg}{ChatGPT 时刻：alignment 改行为，pretraining loss 承载基础能力}{Stanford Online 官方录像 00:13:05}

交叉熵与 perplexity 的关系是
\[
H=-\frac{1}{T}\sum_{t=1}^{T}\log p_\theta(x_t\mid x_{<t}),
\qquad
\operatorname{PPL}=e^H.
\]
perplexity 可直观理解为“模型每个 token 面前的有效选择数”，但不同 tokenizer、语料和 domain 的 PPL 不可直接横比。

\begin{warningbox}{“同 loss 同能力”不是普遍定理}
课堂提出的是强经验判断。即使平均 pretraining loss 相近，architecture、tokenizer、data mixture、optimization trajectory、context length 与 evaluation format 仍可造成能力差异。2024 年 9 月之后的工作也给出 same-loss、different-downstream 的反例，因此本讲只保留其预算与 over-training 直觉。
\end{warningbox}

\teachervoice{00:10:00--00:13:18，讲者明确把所谓 emergent ability 重新解释为 loss 曲线现象。这是他的研究立场，不应改写成社区已经证明的因果结论。}

\subsection{本章小结}

BERT moment 关注训练目标，GPT-3 moment 关注规模与预算，ChatGPT moment 关注低成本行为适配。三者共同把研究问题从“哪种任务最聪明”转成“知识由何种数据和计算获得、行为又怎样被高质量反馈塑形”。

